\chapter{El TypeGraph de sotFS}
\label{cap:typegraph}

\emph{Antes de hablar de operaciones, hace falta fijar la ontología sobre la
que esas operaciones reescriben.} Este capítulo define el \emph{TypeGraph}
de sotFS: la estructura de datos sobre la que el filesystem entero se
construye. Después de leerlo, el lector tiene las definiciones formales de
los seis tipos de nodo, las siete aristas tipadas y los siete invariantes
que el runtime verifica tras cada reescritura.

Este capítulo y el capítulo~\ref{cap:dpo} son los más densos del libro: si
algún día hay que reabrir un único capítulo para entender qué objetos viven
en el grafo, es éste.

\section{Por qué un grafo tipado y no un árbol}

Un filesystem POSIX clásico (\texttt{ext4}, \texttt{XFS}, \texttt{NTFS})
modela su metadata como un \emph{árbol de inodes}: los directorios son
nodos internos, los archivos son hojas, y la única manera de \emph{tener
dos nombres} para un inode es el bolt-on llamado \emph{hard link}, que
rompe la disciplina arbórea.

sotFS abandona el árbol y modela todo como un \emph{grafo dirigido tipado}.
Más concretamente:

\begin{itemize}
  \item Los \emph{vértices} están particionados en seis tipos: inodes,
    directorios, capabilities, transacciones, versiones y bloques. Cada
    tipo tiene su propio conjunto de atributos.
  \item Las \emph{aristas} están particionadas en siete tipos
    (\econtains, \egrants, \edelegates, \ederivedfrom, \esupersedes,
    \epointsto, \ehasxattr); cada arista lleva una etiqueta o atributos
    propios (un nombre de archivo, un bitmask de derechos, un offset).
  \item Cada operación POSIX se ejecuta como una \emph{reescritura} sobre
    el grafo, sujeta a un conjunto fijo de \emph{invariantes
    estructurales} que el sistema verifica después de cada paso.
\end{itemize}

La elección no es estética. Permite codificar como \emph{aristas tipadas}
mecanismos que el árbol clásico tiene que codificar como \emph{convenciones
fuera de banda}: hard links, snapshots, capabilities, dependencias de
transacciones. Cuando una propiedad de seguridad o de consistencia se
puede expresar como ``ningún ciclo en cierta sub-relación'' o ``el
\emph{rights} del hijo es subconjunto del del padre'', tener la relación
como una arista en un grafo tipado significa que el mecanismo de
verificación es uniforme: el chequeador de invariantes recorre el grafo,
no inspecciona estructuras heterogéneas.

\section{Definición formal del TypeGraph}

\begin{defn}[TypeGraph]
\label{def:typegraph}
Un \emph{TypeGraph} es una séxtupla
\[
  G \;=\; \bigl(V,\; E,\; \mathrm{src},\; \mathrm{tgt},\; \tau_V,\; \tau_E,\; \mathrm{attr}\bigr)
\]
donde:
\begin{itemize}
  \item $V$ es un conjunto finito de vértices.
  \item $E$ es un conjunto finito de aristas dirigidas.
  \item $\mathrm{src}: E \to V$ y $\mathrm{tgt}: E \to V$ asignan extremos a
    cada arista.
  \item $\tau_V: V \to T_V$ asigna a cada vértice uno de los seis tipos
    \[
      T_V \;=\; \{\NInode,\; \NDir,\; \NCap,\; \NTxn,\; \NVer,\; \NBlk\}.
    \]
  \item $\tau_E: E \to T_E$ asigna a cada arista uno de los siete tipos
    \[
      T_E \;=\; \{\econtains,\; \egrants,\; \edelegates,\; \ederivedfrom,\; \esupersedes,\; \epointsto,\; \ehasxattr\}.
    \]
  \item $\mathrm{attr}$ asigna a cada vértice y cada arista los atributos
    propios de su tipo (definidos en las secciones~\ref{sec:nodos}
    y~\ref{sec:aristas}).
\end{itemize}
Los pares $(\tau_E, (\tau_V, \tau_V))$ admisibles están restringidos por
una \emph{signatura} fija $\sigma$ (Tabla~\ref{tab:signature}): no toda
arista puede unir cualquier par de tipos.
\end{defn}

\begin{observacion}[Implementación en el runtime]
La séxtupla anterior se materializa en Rust como la \texttt{struct}
\texttt{TypeGraph} definida en \citaLine{sotfs-graph/src/graph.rs}{52}. Los
conjuntos $V$ se almacenan en \emph{arenas} tipadas (una por tipo de
vértice) con \emph{free-list} de slots reusables; el conjunto $E$ vive en
una sola arena de aristas. Los índices auxiliares (\texttt{dir\_contains},
\texttt{cap\_parent}, \texttt{dir\_name\_idx}, \dots) son derivados, no
parte del estado canónico, y se reconstruyen tras una deserialización
(cf.\ cap.~\ref{cap:persistencia}).
\end{observacion}

\section{Los seis tipos de nodo}
\label{sec:nodos}

\begin{defn}[Inode, Directory, Capability, Transaction, Version, Block]
\label{def:nodos}
Cada tipo de nodo tiene los atributos de la
Tabla~\ref{tab:nodos}.
\end{defn}

\begin{table}[h]
\centering
\renewcommand{\arraystretch}{1.15}
\begin{tabularx}{\linewidth}{l X}
  \toprule
  \textbf{Tipo} & \textbf{Atributos} \\
  \midrule
  \NInode & \texttt{id, vtype, permissions, uid, gid, size, link\_count,
                   ctime, mtime, atime} \\
  \NDir & \texttt{id, inode\_id} (el inodo que ``.'' apunta) \\
  \NCap & \texttt{id, rights} (bitmask u8), \texttt{epoch} (generación) \\
  \NTxn & \texttt{id, tier, state, read\_set, write\_set, begin\_ts} \\
  \NVer & \texttt{id, timestamp, root\_inode\_id} \\
  \NBlk & \texttt{id, sector\_start, sector\_count, refcount} \\
  \bottomrule
\end{tabularx}
\caption{Atributos de cada tipo de nodo. Definidos en
\citaLine{sotfs-graph/src/types.rs}{91} (Inode), \texttt{:142} (Directory),
\texttt{:174} (Capability), \texttt{:227} (Transaction), \texttt{:238}
(Version), \texttt{:246} (Block).}
\label{tab:nodos}
\end{table}

\subsection*{Notas sobre cada tipo}

\paragraph{\NInode.} Es la metadata POSIX clásica. Notar que el contenido
de archivos no vive en el inode: vive en \emph{nodos} \NBlk{} a los que el
inode apunta vía aristas \epointsto. \texttt{link\_count} es la cuenta de
hard links, sostenida como invariante por el sistema (cf.\
sec.~\ref{sec:invariantes}).

\paragraph{\NDir.} Lleva sólo dos campos: su \texttt{id} y el
\texttt{inode\_id} del inodo que el ``.'' del directorio apunta. La
metadata del directorio (permisos, owner) vive en su \NInode{} pareado, no
en el \NDir{} mismo. Esta separación es deliberada: el \NDir{} es el
\emph{namespace} (un conjunto de aristas \econtains{} salientes), y el
\NInode{} es la \emph{metadata}.

\paragraph{\NCap.} Una \emph{capability} es un nodo de primera clase, no
una bolt-on de POSIX. Cada cap tiene un \texttt{rights} de 8 bits
(\texttt{READ}, \texttt{WRITE}, \texttt{EXECUTE}, \texttt{GRANT},
\texttt{REVOKE}) y un \texttt{epoch} para invalidación. El cap.~\ref{cap:capabilities}
desarrolla el modelo completo.

\paragraph{\NTxn.} Una transacción es un nodo. Tiene estado
(\texttt{Active}, \texttt{Preparing}, \texttt{Committed}, \texttt{Aborted}),
\emph{tier} (Tier~0/1/2 mapean a las jerarquías de SOT) y dos conjuntos
\texttt{read\_set}/\texttt{write\_set} usados para la detección de
conflictos.

\paragraph{\NVer.} Una versión es un \emph{snapshot pointer}: un timestamp
y un \texttt{root\_inode\_id}. La cadena \ederivedfrom{} entre versiones
forma el linaje de snapshots; es DAG, no árbol.

\paragraph{\NBlk.} Un bloque de datos es un \emph{extent} con
\texttt{sector\_start}, \texttt{sector\_count} y \texttt{refcount}. El
refcount es la cuenta de inodes que lo apuntan vía \epointsto; es
crítico para el GC del backend de almacenamiento
(cap.~\ref{cap:persistencia}).

\section{Los siete tipos de arista}
\label{sec:aristas}

\begin{defn}[Aristas tipadas y signatura]
\label{def:aristas}
Cada arista pertenece a uno de siete tipos. Su signatura
$\tau_E \to \tau_V \times \tau_V$ y atributos propios están en la
Tabla~\ref{tab:signature}.
\end{defn}

\begin{table}[h]
\centering
\renewcommand{\arraystretch}{1.2}
\begin{tabularx}{\linewidth}{l l l X}
  \toprule
  \textbf{Tipo} & \textbf{Origen} & \textbf{Destino} & \textbf{Atributos / sentido} \\
  \midrule
  \econtains      & \NDir & \NInode    & \texttt{name}: nombre POSIX bajo el que el inodo aparece en el dir.\\
  \egrants        & \NCap & \NInode    & \texttt{rights}: subconjunto de los rights de la cap.\\
  \edelegates     & \NCap & \NCap      & forma el árbol CDT (\emph{capability derivation tree}).\\
  \ederivedfrom   & \NVer & \NVer      & linaje de snapshots; DAG.\\
  \esupersedes    & \NInode & \NInode  & reemplazo atómico durante \oprename.\\
  \epointsto      & \NInode & \NBlk    & \texttt{offset}: posición del bloque dentro del archivo.\\
  \ehasxattr      & \NInode & \NInode (\textsf{XAttr}) & cada xattr es un nodo de primera clase.\\
  \bottomrule
\end{tabularx}
\caption{Signatura de las siete aristas. La columna \emph{Origen} y
\emph{Destino} fijan los tipos admisibles; un par fuera de la signatura
viola el TypeInvariant. El \textsf{XAttr} aparece como séptimo tipo de
nodo en algunas vistas; en este libro lo tratamos como una variante de
\NInode{} para mantener el conteo en seis. La definición canónica vive en
\citaLine{sotfs-graph/src/types.rs}{346}.}
\label{tab:signature}
\end{table}

\subsection*{Por qué siete y no menos}

Cada arista codifica una relación que en un FS POSIX clásico está
\emph{disuelta} en convenciones, en estructuras separadas o en
\emph{flags}. Cuando se empuja la relación a una arista tipada, dos cosas
mejoran:

\begin{itemize}
  \item La verificación es uniforme. ``\textsf{No hay ciclos vía
    \edelegates}'' es el mismo tipo de propiedad que ``\textsf{No hay
    ciclos vía \econtains{} excluyendo \texttt{..}}''. Ambas se chequean
    con el mismo molde.
  \item La extensión es modular. Un futuro tipo de arista
    (digamos \textsf{auditedBy}: \NInode{} a un nodo de auditoría) se
    suma a la signatura sin reescribir todo.
\end{itemize}

La trampa simétrica es que cada arista nueva amplía la superficie de
invariantes que hay que probar y mantener. El cap.~\ref{cap:dpo} muestra
que cada operación POSIX se diseña explícitamente para preservar todas
las invariantes; el cap.~\ref{cap:coq} muestra hasta dónde esa
preservación está mecanizada hoy y dónde quedan deudas (los cinco
\texttt{Admitted}).

\begin{observacion}[Evolución del modelo]
\label{obs:evolucion-aristas}
La versión preliminar de sotFS (la que vivía como subsistema de sotX,
caps.\ 17--18 de su libro) listaba seis aristas:
\econtains, \texttt{extent}, \texttt{named-as}, \texttt{parent},
\texttt{captured} y \texttt{owns}. La reorganización de la versión
\texttt{v0.2.0} colapsó algunas de esas relaciones y agregó otras: la
arista \texttt{extent} pasó a ser \epointsto{} con \texttt{offset};
\texttt{named-as} se fundió con \econtains{} (el nombre vive como
atributo de la arista, no como nodo intermedio); \texttt{parent} se
suprimió y se deriva por consulta inversa con el índice
\texttt{dir\_for\_inode\_idx}; \texttt{captured} se generalizó a
\egrants; \texttt{owns} se reemplazó por el mecanismo de \texttt{domain\_id}
en el \texttt{CapContext}. Las nuevas \edelegates, \ederivedfrom,
\esupersedes{} y \ehasxattr{} no existían como aristas: estaban
codificadas como bolt-ons (snapshots como copias, xattrs como
diccionarios laterales). El libro adopta el modelo del runtime
v0.2.5; este recuadro existe sólo para evitar confusión a quien venga
con la versión vieja en mano.
\end{observacion}

\section{Diagrama de un TypeGraph mínimo}

La figura~\ref{fig:tg-mini} muestra un TypeGraph well-formed con un
directorio raíz, un subdirectorio, dos archivos que comparten un bloque
de datos, y una capability que autoriza acceso a uno de los inodes.

\begin{figure}[h]
\centering
\begin{tikzpicture}[node distance=10mm and 14mm]
  % nodos
  \node[ndir] (root)      {Dir\\\texttt{D1}};
  \node[nino, above=4mm of root] (rooti) {Inode\\\texttt{I1}};
  \node[ndir, right=20mm of root] (sub) {Dir\\\texttt{D2}};
  \node[nino, above=4mm of sub] (subi) {Inode\\\texttt{I2}};
  \node[nino, below right=14mm and 6mm of sub] (i3) {Inode\\\texttt{I3}};
  \node[nino, below=18mm of sub] (i4) {Inode\\\texttt{I4}};
  \node[nblk, below=10mm of i3] (b1) {Block\\\texttt{B1}};
  \node[nblk, below=10mm of i4] (b2) {Block\\\texttt{B2 (shared)}};
  \node[ncap, left=22mm of i4] (c1) {Cap\\\texttt{C1}\\rights=R/W};

  % contains: dir → inode
  \draw[arr] (root) -- (rooti) node[edgelbl,midway] {.};
  \draw[arr] (root) -- (sub) node[edgelbl,midway] {sub};
  \draw[arr] (sub) -- (subi) node[edgelbl,midway] {.};
  \draw[arr] (sub) to[bend left=10] node[edgelbl, midway] {f1.txt} (i3);
  \draw[arr] (sub) to[bend right=10] node[edgelbl, near end] {f2.txt} (i4);

  % pointsTo: inode → block
  \draw[arr-points] (i3) -- (b1) node[edgelbl, midway] {pointsTo};
  \draw[arr-points] (i3) -- (b2) node[edgelbl, near start, sloped] {pointsTo};
  \draw[arr-points] (i4) -- (b2) node[edgelbl, midway] {pointsTo};

  % grants
  \draw[arr-grant] (c1) -- (i3) node[edgelbl, midway, sloped] {grants R};
\end{tikzpicture}
\caption{Un TypeGraph well-formed mínimo con dos directorios, dos
archivos compartiendo el bloque \texttt{B2}, y una cap autorizando lectura
de \texttt{I3}. Las aristas \econtains{} son sólidas (negro); las
\epointsto{} verdes; las \egrants{} violetas dashed.}
\label{fig:tg-mini}
\end{figure}

Notar tres cosas que el árbol POSIX no expresaría así:

\begin{enumerate}
  \item \texttt{B2} es un único nodo apuntado por dos inodes distintos;
    en \texttt{ext4} esa situación requeriría duplicar el bloque (no hay
    sharing nativo) o hacer reflinks (extensión específica de XFS/btrfs).
  \item La cap \texttt{C1} es un nodo, no una entrada en una tabla
    paralela. La autorización a \texttt{I3} es la arista \egrants.
  \item Cada \NDir{} tiene un \NInode{} pareado (el \texttt{inode\_id} del
    \NDir, alcanzable por la arista ``\,.\,''). Esa pareja es lo que el
    \invref{inv:dir-self} obliga a mantener.
\end{enumerate}

\section{Los siete invariantes verificados tras cada paso}
\label{sec:invariantes}

El runtime expone una función \texttt{check\_invariants()}
(\citaLine{sotfs-graph/src/graph.rs}{874}) que se llama tras cada paso DPO
en el camino de tests, en los fuzzers y, opcionalmente, en producción
(modo paranoid). Verifica siete invariantes más una consistencia de índice
secundario.

\begin{invariante}[Link-count consistency, I2 + G3]
\label{inv:link-count}
Para todo $i \in \NInode$,
\[
  i.\texttt{link\_count}
  \;=\;
  \bigl|\{\, e \in E : \tau_E(e) = \econtains \;\wedge\; \mathrm{tgt}(e) = i \;\wedge\; e.\texttt{name} \neq \texttt{".."} \,\}\bigr|.
\]
Es decir, la cuenta de hard links almacenada en el inode coincide con la
cantidad de aristas \econtains{} entrantes, excluyendo las aristas
``..''. Implementado en
\citaLine{sotfs-graph/src/graph.rs}{888}.
\end{invariante}

\begin{invariante}[Unique names per directory, C1 + I4]
\label{inv:unique-names}
Para todo $d \in \NDir$ y todo par de aristas $e_1, e_2$ \econtains{} que
salen de $d$,
\[
  e_1 \neq e_2 \;\Longrightarrow\; e_1.\texttt{name} \neq e_2.\texttt{name}.
\]
Implementado en \citaLine{sotfs-graph/src/graph.rs}{918}.
\end{invariante}

\begin{invariante}[Dir self-reference, I3]
\label{inv:dir-self}
Para todo directorio $d \in \NDir$ existe una arista \econtains{} desde
$d$ con \texttt{name = "."} y destino $d.\texttt{inode\_id}$. Es decir,
la pareja $\NDir/\NInode$ siempre se cierra sobre sí misma vía la entrada
``.''. Implementado en \citaLine{sotfs-graph/src/graph.rs}{936}.
\end{invariante}

\begin{invariante}[No dangling edges, G2]
\label{inv:no-dangling}
Para toda arista $e \in E$, $\mathrm{src}(e), \mathrm{tgt}(e) \in V$:
ambos extremos son nodos vivos. La consecuencia operativa es que la
remoción de un nodo obliga a remover \emph{todas} sus aristas adyacentes
en el mismo paso DPO. Implementado en
\citaLine{sotfs-graph/src/graph.rs}{964}.
\end{invariante}

\begin{invariante}[Block refcount, B1]
\label{inv:block-refcount}
Para todo bloque $b \in \NBlk$,
\[
  b.\texttt{refcount}
  \;=\;
  \bigl|\{\, e \in E : \tau_E(e) = \epointsto \;\wedge\; \mathrm{tgt}(e) = b \,\}\bigr|.
\]
Vital para la corrección del GC en el backend de almacenamiento: si el
refcount divergiera del conjunto real de aristas \epointsto, el GC
liberaría bloques aún apuntados. Implementado en
\citaLine{sotfs-graph/src/graph.rs}{996}.
\end{invariante}

\begin{invariante}[No directory cycles, A1]
\label{inv:no-dir-cycles}
La sub-relación inducida por \econtains{} restringida a directorios
(ignorando ``\,.\,'' y ``\,..\,'') es acíclica: el conjunto de
directorios alcanzables desde cualquier directorio dado no se vuelve
sobre sí mismo. Implementado en \citaLine{sotfs-graph/src/graph.rs}{1015}
con DFS desde cada raíz; cf.\ la trampa~\ref{tr:on-cube}.
\end{invariante}

\begin{invariante}[Capability monotonicity, G4]
\label{inv:cap-monotonicity}
Para todo par de capabilities $c, c' \in \NCap$ unidas por una arista
\edelegates{} de $c$ a $c'$,
\[
  c'.\texttt{rights} \;\subseteq\; c.\texttt{rights}.
\]
Una delegación nunca crea derechos; sólo restringe los del padre.
Implementado en \citaLine{sotfs-graph/src/graph.rs}{1056}.
\end{invariante}

A las anteriores se suma la consistencia del índice secundario
\texttt{dir\_name\_idx} (\citaLine{sotfs-graph/src/graph.rs}{528}), que
no es invariante \emph{semántico} sino correspondencia entre dos
representaciones del mismo dato y se chequea aparte.

\subsection*{Algoritmo de verificación}

El \texttt{check\_invariants()} ejecuta los siete chequeos en cascada y
aborta al primer fallo, devolviendo un \texttt{GraphError::InvariantViolation}.

\begin{algorithm}[H]
\SetAlgoLined
\KwIn{TypeGraph $G$}
\KwOut{\textbf{ok} si $G$ es well-formed; error si no.}
\caption{\textsc{CheckInvariants}}\label{alg:check-invariants}

\textsc{CheckLinkCountConsistency}($G$)\;
\textsc{CheckUniqueNames}($G$)\;
\textsc{CheckDirSelfRef}($G$)\;
\textsc{CheckNoDanglingEdges}($G$)\;
\textsc{CheckBlockRefcount}($G$)\;
\textsc{CheckNoDirCycles}($G$)\;
\textsc{CheckCapMonotonicity}($G$)\;
\textsc{CheckDirNameIdxConsistency}($G$)\;
\Return{\textbf{ok}}\;
\end{algorithm}

\subsection*{Costo}

Sea $N = |V|$ y $M = |E|$. La Tabla~\ref{tab:cost-invariants} resume el
costo de cada chequeo.

\begin{table}[h]
\centering
\begin{tabular}{lll}
  \toprule
  \textbf{Chequeo} & \textbf{Costo} & \textbf{Notas} \\
  \midrule
  link-count   & $O(N + M)$ & barrido lineal de inodes y \econtains{}. \\
  unique-names & $O(M \log M)$ & \texttt{BTreeSet} por directorio. \\
  dir-self-ref & $O(N \log M)$ & lookup por dir. \\
  no-dangling  & $O(M)$ & barrido de aristas. \\
  block-refcount & $O(N + M)$ & barrido de bloques y \epointsto. \\
  no-dir-cycles & $O(N \cdot \mathrm{depth})$ típico, $O(N^3)$ worst-case$^\dagger$ & DFS por dir. \\
  cap-monotonicity & $O(|\NCap|)$ & barrido de aristas \edelegates. \\
  dir-name-idx & $O(M \log N)$ & comparar índice contra \econtains. \\
  \bottomrule
\end{tabular}
\caption{Costo de cada chequeo. $^\dagger$ El worst-case $O(N^3)$ del
chequeo de ciclos viene de un patrón patológico de cadenas de
\opmkdir{} profundas; está documentado y mitigable con un índice
adicional. Cf.\ trampa~\ref{tr:cycle-cost}.}
\label{tab:cost-invariants}
\end{table}

\begin{trampa}[Costo cúbico del cycle check con \texttt{deep\_mkdir\_chain}]
\label{tr:cycle-cost}
\citaCommit{(M4.1.1)}. Versión vulnerable: \texttt{dir\_for\_inode}
hacía un barrido lineal sobre \texttt{self.dirs.values()} para encontrar
el \NDir{} cuyo \texttt{inode\_id} coincidiera con el dado. La función se
llama dentro de \texttt{is\_descendant\_of} y \texttt{has\_cycle\_from},
una vez por arista durante un DFS. El loop externo del chequeo
\textsc{CheckNoDirCycles} itera todos los dirs y arranca un DFS desde
cada uno: el costo se vuelve $O(N^3)$ en peor caso. La proptest
\texttt{deep\_mkdir\_chain\_no\_cycles} construye cadenas de 20--30
\opmkdir{} anidados; con 5 cases corre en 0{,}01 s, con 50 cases pasa de
60 s.

\textbf{Fix} (M4.1.1): un índice \texttt{dir\_for\_inode\_idx:
BTreeMap<InodeId, DirId>} mantenido junto con \texttt{insert\_dir} y
\texttt{remove\_dir} reduce \texttt{dir\_for\_inode} a $O(\log N)$ y el
chequeo entero a $O(N \cdot \mathrm{depth})$ típico
(\citaLine{sotfs-graph/src/graph.rs}{109}).

\textbf{Lección.} Una invariante puede ser \emph{correcta} y el chequeo
\emph{insostenible} bajo workload patológico. La proptest sirvió para
detectar el problema antes de que produciera un timeout en CI; el fix
fue sumar un índice secundario, no relajar la propiedad.
\end{trampa}

\begin{trampa}[\texttt{lookup\_name} cuadrático bajo fill hostil]
\label{tr:on-cube}
\citaCommit{(v0.2.0 → v0.2.4)}. Versión vulnerable:
\texttt{lookup\_name(dir, name)} barría linealmente
\texttt{dir\_contains[dir]} comparando el \texttt{name} de cada arista.
Bajo un workload donde un usuario llena un directorio compartido con $N$
nombres, cada \texttt{stat}/\texttt{lookup} contra ese dir era $O(N)$;
\texttt{perf record} mostró \texttt{lookup\_name + memcmp} a $\sim$92\%
del CPU del daemon FUSE. Vector de DoS local trivial.

\textbf{Fix.} \texttt{dir\_name\_idx: BTreeMap<(DirId, String), EdgeId>},
mantenido por \texttt{insert\_edge}, \texttt{remove\_edge} y
\texttt{rename\_contains\_edge}; \texttt{lookup\_name} pasa a $O(\log N)$
(\citaLine{sotfs-graph/src/graph.rs}{111}).

\textbf{Lección.} Las propiedades estructurales no se imponen solas en
runtime: hay que diseñar las índices que las hagan baratas. La
verificación de la propiedad ``\,$\texttt{stat}$ es flat en tamaño de
directorio'' es indirecta: el invariante semántico (unicidad de
nombres) ya estaba probado, pero el costo de chequearlo o de aprovecharlo
para una lookup era el cuello de botella.
\end{trampa}

\section{La cara TLA\textsuperscript{+} del TypeGraph}

El \citaTLA{sotfs\_graph.tla} (518 líneas) modela el TypeGraph como una
máquina de estados sobre los mismos seis tipos de nodo. El estado se
expresa con variables como las siguientes:

\begin{lstlisting}[language=tla, caption={Extracto del estado en
\texttt{sotfs\_graph.tla}.}, label={lst:tla-state}]
VARIABLES
    inodes,        \* Set of active InodeIds
    dirs,          \* Set of active DirIds
    caps,          \* Set of active CapIds
    blocks,        \* Set of active BlockIds
    contains,      \* Set of <<DirId, Name, InodeId>> triples
    grants,        \* Function: CapId -> InodeId
    delegates,     \* Function: CapId -> CapId (parent)
    rights         \* Function: CapId -> SUBSET RightsAlphabet
\end{lstlisting}

Las acciones (\texttt{Mkdir}, \texttt{Unlink}, \texttt{Link},
\texttt{Rename}, \dots) se enuncian con su \emph{gluing condition} a la
izquierda del \texttt{/\textbackslash} y la actualización del estado a la
derecha. El invariante \texttt{TypeOK} restringe los tipos de los
elementos en cada conjunto; \texttt{NameUniqueness},
\texttt{LinkCountConsistent} y \texttt{NoDirCycles} son los análogos
formales de los invariantes~\ref{inv:link-count},
\ref{inv:unique-names} y \ref{inv:no-dir-cycles}.

El cap.~\ref{cap:tla} cubre cómo correr TLC sobre el spec, qué configs
small/medium/large se usan y cómo leer un counterexample.

\section{Resumen del capítulo}

\begin{itemize}
  \item Un \emph{TypeGraph} es una séxtupla
    $(V, E, \mathrm{src}, \mathrm{tgt}, \tau_V, \tau_E, \mathrm{attr})$
    con seis tipos de nodo y siete tipos de arista restringidos por una
    signatura fija.
  \item Cada nodo tiene atributos POSIX-friendly (Tabla~\ref{tab:nodos});
    cada arista lleva un atributo propio (nombre, rights, offset).
  \item Siete invariantes (link-count, unique-names, dir-self-ref,
    no-dangling, block-refcount, no-dir-cycles, cap-monotonicity) más
    la consistencia del índice \texttt{dir\_name\_idx} se verifican tras
    cada paso DPO. Su algoritmo es \textsc{CheckInvariants}
    (Alg.~\ref{alg:check-invariants}).
  \item Dos trampas históricas: el chequeo de ciclos en peor caso
    $O(N^3)$ (mitigado con \texttt{dir\_for\_inode\_idx}) y el
    \texttt{lookup\_name} lineal (mitigado con \texttt{dir\_name\_idx}).
    Ambos invariantes \emph{semánticos} eran correctos antes; los fixes
    fueron \emph{índices secundarios} para hacer el chequeo y la
    operación baratos en runtime.
  \item El TypeGraph del runtime y el del \citaTLA{sotfs\_graph.tla} se
    corresponden por construcción: ambos manipulan los mismos seis
    tipos y las mismas siete relaciones; la verificación TLC explora un
    espacio de estados acotado y certifica que las acciones preservan
    los invariantes (cap.~\ref{cap:tla}).
\end{itemize}

\section*{Ejercicios}

\begin{ejercicio}[\dificultad{$\star$}]
\label{ej:tg-1}
Enuncie los siete invariantes de la sección~\ref{sec:invariantes} en una
línea cada uno. ¿Cuál depende de un invariante \emph{numérico} (cuenta)
y no estructural?
\end{ejercicio}

\begin{ejercicio}[\dificultad{$\star$}]
\label{ej:tg-2}
Liste los seis tipos de nodo y los siete tipos de arista. Para cada
arista, indique los tipos de origen y destino admisibles según la
signatura $\sigma$ (Tabla~\ref{tab:signature}).
\end{ejercicio}

\begin{ejercicio}[\dificultad{$\star\star$}]
\label{ej:tg-3}
Dibuje el TypeGraph después de las operaciones \texttt{mkdir /a};
\texttt{create /a/x.txt}; \texttt{create /a/y.txt}; \texttt{ln /a/x.txt
/a/z.txt}. ¿Cuál es el \texttt{link\_count} del inodo de \texttt{x.txt}
después de la última operación? ¿Cuántas aristas \econtains{} entran al
inodo de \texttt{x.txt} (excluyendo ``..'')?
\end{ejercicio}

\begin{ejercicio}[\dificultad{$\star\star$}]
\label{ej:tg-4}
Demuestre que el \invref{inv:link-count} \emph{implica} que un inode
con \texttt{link\_count = 0} no tiene aristas \econtains{} entrantes
(salvo ``..''). Use solo la definición.
\end{ejercicio}

\begin{ejercicio}[\dificultad{$\star\star$}]
\label{ej:tg-5}
Considere un TypeGraph con dos directorios $D_1, D_2$ donde
$D_1.\texttt{inode\_id} = i_1$ y $D_2.\texttt{inode\_id} = i_2$. Si una
arista \econtains{} con \texttt{name="."} sale de $D_1$ y apunta a
$i_2$ ($\neq i_1$), ¿cuál invariante se rompe? ¿En qué línea de
\texttt{check\_invariants} se detecta?
\end{ejercicio}

\begin{ejercicio}[\dificultad{$\star\star$}]
\label{ej:tg-6}
La trampa~\ref{tr:cycle-cost} se cierra agregando
\texttt{dir\_for\_inode\_idx}. Discuta el costo de mantener ese índice
\emph{coherente} a lo largo de \texttt{insert\_dir} y \texttt{remove\_dir}.
¿Bajo qué patrón de operaciones es más caro mantener el índice que
recorrer linealmente?
\end{ejercicio}

\begin{ejercicio}[\dificultad{$\star\star$}]
\label{ej:tg-7}
Una arista \edelegates{} de $c$ a $c'$ con $c'.\texttt{rights} \not\subseteq
c.\texttt{rights}$ viola el \invref{inv:cap-monotonicity}. Diseñe el
test unitario más corto posible que detecte la violación: ¿qué grafo de
tres nodos hace falta? ¿qué llamada al runtime?
\end{ejercicio}

\begin{ejercicio}[\dificultad{$\star\star\star$}]
\label{ej:tg-8}
Diseñe un octavo invariante para detectar \emph{leaks de bloques}: un
\NBlk{} cuyo \texttt{refcount > 0} pero cuyo conjunto de aristas
\epointsto{} entrantes está vacío (caso patológico de divergencia entre
contador y aristas). Justifique cuándo lo correrías (¿en cada paso DPO?
¿solo en el GC?). Pista: pensá la asimetría entre el costo de
mantener y el de chequear.
\end{ejercicio}

\begin{ejercicio}[\dificultad{$\star\star\star$}]
\label{ej:tg-9}
Extienda la signatura $\sigma$ con un octavo tipo de arista
\textsf{auditedBy} (\NInode{} a un nuevo nodo \textsf{Audit}) que registre
quién auditó la última escritura. Discuta: ¿qué invariante hay que
agregar? ¿qué operaciones POSIX hay que modificar para mantenerlo? ¿qué
spec TLA\textsuperscript{+} hay que extender?
\end{ejercicio}
